\section{Limitations}

Our experiment faces several limitations that might weaken our claims. 

\textbf{Data}

The questionnaire responses were anonymized, not allowing us to connect students' behavior and performance to their attitudes. 

The groups in our experiment were not of same size, with some participants dropping out during the course (at the end of the course: \ac{pv} - 21, \ac{cv} - 14). That might affect findings from comparing the introductory and feedback questionnaires. The relatively small sample size reduces the statistical power of our findings.

On top of that, the participants did not have the same baseline knowledge across both courses, making the interpretation of collected data more complicated.

Lastly, it should be noted that the difference in overall satisfaction with the course as well as interest in e-learning in general might be at least partially attributed to a potentially better design of \ac{pv} compared to \ac{cv}. However, there is no way to verify the assumption without further experiments.


\textbf{Course design}

One hour long course does not allow for implementation of some concepts, such as techniques related to gradual increase in difficulty, meaningful progressions, or advanced adaptive learning techniques. Additionally, an activity that spans a longer time frame might display more significant differences in observed groups.

% Another drawback was an inability to engage students in real-time group activities,
Another drawback was the combination of the absence of an automated diagram assessment tool and the unavailability of teacher intervention. This led us to substitute direct modeling exercises with alternatives that were easier to correct automatically.


% \begin{itemize}
    % \item Anonymized questionnaires
    % \item too few participants and uneven start skills and not the same number of participants across both courses
    % \item some data gathered is tricky to interpret - time, interactions \todo{any other issues with data?}
    % \item size and time of the course - not allow for some of the more advanced activities in e-learning, cannot overload the students
    % \item baseline kknowledge - students were already familiar with the class diagram
    % \item setting - impossibility of direct group activities
    % \item exercises - would be better to have an automated assessment tool and use draw.io and some diagram-code conversion tool instead of adapting general purpose exercise types, but outside the scope of 1 DP project
    % \item moodle - what is difficult to do there?
    % \item activities restricted to "no intervention from teachers" - no larger or more complex exercises, only automated correction
    % \item visual animations - we do not have that now, maybe leaderboard like in kahoot and such, transitions, changes based on performance...
    % \item nepytal som sa v dotazniku na to, ake aktivity sa im pacili, to som ale kar ja a nemal by som poukazovat na neobhajitelne zle dizajn rozhodnutie
% \end{itemize}